\documentclass[12pt,a4paper]{article}

\usepackage[spanish]{babel}
\usepackage[utf8]{inputenc}
\usepackage[T1]{fontenc}
\usepackage{geometry}
\usepackage{amsmath,amssymb,amsthm}
\usepackage{graphicx}
\usepackage{hyperref}

\geometry{margin=2.5cm}

\title{Exposición 1: Igualdad de Parseval}
\author{Lester Iván Hernández Hernández (20242000637)\\
        Emilio Gabriel Reyes Pérez (20242000121)}
\date{\today}

\begin{document}

\maketitle

\section{Enunciado del problema}

Sea $f(x) \in L^2[-p,p] \longrightarrow \mathbb{R}$ cuyo desarrollo de Fourier está dado por
\[
f(x) = \frac{a_0}{2} + \sum_{n=1}^{\infty} a_n \cos \frac{n\pi}{p}x + b_n \operatorname{sen} \frac{n\pi}{p}x
\]

Probar la igualdad de Parseval
\[
\langle f,f\rangle = \|f\|^2 = \int_{-p}^{p} (f(x))^2\,dx = p\left(\frac{a_0^2}{2} + \sum_{n=1}^{\infty} (a_n^2 + b_n^2)\right)
\]

\section{Definición de L2}
Una sucesión de funciones $\{f_n\}$ converge a una función $f$ en la norma $L^2$ si la integral de la diferencia al cuadrado tiende a cero cuando $n$ tiende a infinito, es decir,
\[
\lim_{n \to \infty} \int_{-p}^{p} |f_n(x) - f(x)|^2 \, dx = 0
\]

\section{Solución}

Siguiendo la definción de norma con función peso $1$, 

\begin{align*}
\|f\|^2 &= \int_{-p}^{p} (f(x))^2 \, dx \\
&= \int_{-p}^{p} \left( \frac{a_0}{2} + \sum_{n=1}^{\infty} \left(a_n \cos \frac{n\pi}{p}x + b_n \sin \frac{n\pi}{p}x \right)\right)^2 \, dx 
\end{align*}

Utilizando la identidad del cuadrado de un binomio,

\begin{align*}
\int_{-p}^{p} \left(\frac{a_0^2}{4} + a_0\sum_{n=1}^{\infty} \left(a_n \cos \frac{n\pi}{p}x + b_n \sin \frac{n\pi}{p}x \right) + \left(\sum_{n=1}^{\infty} a_n \cos \frac{n\pi}{p}x + b_n \sin \frac{n\pi}{p}x \right)^2 \right) \, dx
\end{align*}

Ahora, aplicando la linealidad de la integral obtenemos:

\begin{align*}
\int_{-p}^{p} \frac{a_0^2}{4} \, dx + \int_{-p}^{p} a_0\sum_{n=1}^{\infty} \left(a_n \cos \frac{n\pi}{p}x + b_n \sin \frac{n\pi}{p}x \right) \, dx + \int_{-p}^{p} \left(\sum_{n=1}^{\infty} a_n \cos \frac{n\pi}{p}x + b_n \sin \frac{n\pi}{p}x \right)^2 \, dx
\end{align*}

Y ahora resolvemos una por una las integrales:

\begin{align*}
\int_{-p}^{p} \frac{a_0^2}{4} \, dx &= \frac{a_0^2}{4} \int_{-p}^{p} dx = \frac{a_0^2}{4} (2p) = \frac{a_0^2}{2} p
\end{align*}

\begin{align*}
\int_{-p}^{p} a_0\sum_{n=1}^{\infty} \left(a_n \cos \frac{n\pi}{p}x + b_n \sin \frac{n\pi}{p}x \right) \, dx &= a_0 \sum_{n=1}^{\infty} \left( a_n \int_{-p}^{p} \cos \frac{n\pi}{p}x \, dx + b_n \int_{-p}^{p} \sin \frac{n\pi}{p}x \, dx \right) = 0
\end{align*}

\begin{align*}
&\int_{-p}^{p} \left(\sum_{n=1}^{\infty} a_n \cos \frac{n\pi}{p}x + b_n \sin \frac{n\pi}{p}x \right)^2 \, dx 
\\& = \int_{-p}^{p} \left(\sum_{n=1}^{\infty} a_n \cos \frac{n\pi}{p}x + b_n \sin \frac{n\pi}{p}x\right)\left(\sum_{n=1}^{\infty} a_n \cos \frac{n\pi}{p}x + b_n \sin \frac{n\pi}{p}x\right)\, dx
\\& = \int_{-p}^{p} \sum_{n=1}^{\infty} \sum_{m=1}^{\infty} \Bigg(a_n a_m \cos \frac{n\pi}{p}x \cos \frac{m\pi}{p}x + a_n b_m \cos \frac{n\pi}{p}x \sin \frac{m\pi}{p}x + \\&b_n a_m \sin \frac{n\pi}{p}x \cos \frac{m\pi}{p}x + b_n b_m \sin \frac{n\pi}{p}x \sin \frac{m\pi}{p}x\Bigg) \, dx
\\& = \sum_{n=1}^{\infty} \sum_{m=1}^{\infty} \Bigg(a_n a_m \int_{-p}^{p} \cos \frac{n\pi}{p}x \cos \frac{m\pi}{p}x \, dx + a_n b_m \int_{-p}^{p} \cos \frac{n\pi}{p}x \sin \frac{m\pi}{p}x \, dx + \\&b_n a_m \int_{-p}^{p} \sin \frac{n\pi}{p}x \cos \frac{m\pi}{p}x \, dx + b_n b_m \int_{-p}^{p} \sin \frac{n\pi}{p}x \sin \frac{m\pi}{p}x\, dx\Bigg)
\end{align*}

Por ortogonalidad de las funciones seno y coseno, sus integrales se vuelven cero, cuando $n \neq m$. Al igual que las integrales de los productos cruzados, por lo que nos queda:

\begin{align*}
&\sum_{n=1}^{\infty} \sum_{m=1}^{\infty} \Bigg(a_n a_m \int_{-p}^{p} \cos \frac{n\pi}{p}x \cos \frac{m\pi}{p}x \, dx + b_n b_m \int_{-p}^{p} \sin \frac{n\pi}{p}x \sin \frac{m\pi}{p}x\, dx\Bigg) 
\\&= \sum_{n=1}^{\infty} \left(a_n^2 \int_{-p}^{p} \cos^2 \frac{n\pi}{p}x \, dx + b_n^2 \int_{-p}^{p} \sin^2 \frac{n\pi}{p}x\, dx\right)
\\&= \sum_{n=1}^{\infty} \left(a_n^2 p + b_n^2 p\right) 
\\&= p \sum_{n=1}^{\infty} (a_n^2 + b_n^2)
\end{align*}

Por lo cual, se obtiene que:

\begin{align*}
&\int_{-p}^{p} \frac{a_0^2}{4} \, dx + \int_{-p}^{p} a_0\sum_{n=1}^{\infty} \left(a_n \cos \frac{n\pi}{p}x + b_n \sin \frac{n\pi}{p}x \right) \, dx + \int_{-p}^{p} \left(\sum_{n=1}^{\infty} a_n \cos \frac{n\pi}{p}x + b_n \sin \frac{n\pi}{p}x \right)^2 \, dx\\
&= \frac{a_0^2}{2} p + 0 + p \sum_{n=1}^{\infty} (a_n^2 + b_n^2)\\
&= p\left(\frac{a_0^2}{2} + \sum_{n=1}^{\infty} (a_n^2 + b_n^2)\right)
\end{align*}

Que era lo que se quería demostrar, por lo que se concluye que la igualdad de Parseval es verdadera.

\end{document}
